% Section 9 -- Discussion. Author's draft of 2026-09-05, cite keys mapped to refs.bib.
% One figure corrected against the committed data and flagged: the composition of the authored benign material.
\section{Discussion}
\label{sec:discussion}

\textbf{What the minimal-data regime is, and is not.} Our results do not show that activation probes fail; two of
three tasks end with probes performing strongly on real benchmarks after one-line changes to where the activation is
read. Nor do they show the minimal-data regime matches production systems: the frontier deployments we cite train on
data at scales this regime excludes by definition, and DeepMind's report of length-shift failures at production scale
\cite{gdm2026probes} suggests convention sensitivity persists even with that data. What the results show is narrower
and, we think, more useful: in the minimal-data regime, the reported number is a joint property of the probe and of
conventions most papers relegate to implementation detail. Two studies of the same recipe can honestly report
near-perfect and near-chance performance while differing only in template mode, read position, hooked tensor, or the
composition of the benign set. Until those are reported as experimental variables, results in this literature are not
comparable across papers.

\textbf{Implications for practice.} For practitioners, the repairs are free and the checklist is short: read through
the chat template; for localized threats, pool over the relevant span rather than reading the final token; verify
which tensor the serving framework's hook returns; stratify by length; include a topic-matched benign set and a
lexical baseline; separate probe detections from the model's own refusals. For consumers of probe
results---including procurement and deployment decisions---the checklist inverts into due-diligence questions, and a
reported result that cannot answer them should be treated as unmeasured rather than favorable.

\textbf{Limitations.} The obfuscation sweep covers one $9$B model exhaustively rather than many models shallowly; the
finding may not transfer to models that decode more capably. Our benign sets, though pinned and controlled, are partly
authored; Section~\ref{sec:apc}'s own finding cautions against over-trusting authored instruments, and we apply that
caution to ourselves---$40$ of the $120$ agent benign cases are authored, including $26$ of its $36$ hard negatives,
where its positives are real benchmark cases.
Distributional detection is not adversarial robustness: obfuscated-activation attacks \cite{bailey2025obfuscated} show
latent-space monitors can be actively evaded, and nothing here addresses an adaptive attacker. Benchmarks themselves
turn over; the InjecAgent artifacts we document \cite{injecagentissue7} may be repaired, at which point our measured
numbers describe a superseded version. Finally, all probes here are linear; whether the policy-compliance
construct-validity barrier or the encoded-content result changes under non-linear probes is open, though the barrier
argument in Section~\ref{sec:apc} is classifier-independent by construction.

\textbf{Why the protocol section reports its own failures.} Three of our pre-registered designs could not answer their
own questions---a split that leaked injection strings, a comparison whose two arms were algebraically identical, and
an instrument whose validity dissolved as its confounds were controlled. We report these not as confession but as
data: pre-registration's value was not that our designs were right, but that their failures were detectable,
attributable, and correctable in the open. An evaluation regime in which errors that inflate results are as
discoverable as errors that deflate them is the substantive proposal of this paper; the specific conventions are its
current, and certainly incomplete, instantiation.
